\tocsection{\textbf{10/9/26 Lecture 19:} The Fundamental Group of the Circle}

We will show $\pi_1(S^1, p) \simeq \Z$. Heuristically, a loop is sent to its winding number around $0$. Define $\pi : \R \to S^1$ by $\pi(x) = e^{2\pi i x}$.
\fig{0.5}{Lecture19Fig1}{}

Take $1 \in S^1$ as the basepoint. Given $n \in \Z$, say
\[ \gamma_n(s) = ns, \quad 0 \leq s \leq 1 \]
Now, $\pi \circ \gamma_n$ is a loop in $S^1$ that starts and ends at $1$.\\

\textbf{Note:} If $\gamma$ is any path in $\R$ from $0$ to $n$, then $\gamma \simeq \gamma_n$ (rel $\{0, 1\}$) since $\R$ is convex in itself. We will use this result a lot.

\begin{lemma}[\textbf{Lemma 19.1:} Path Lifting Lemma]
    If $\sigma$ is a path in $S^1$ starting at $1$, then there exists a unique path $\tilde{\sigma}$ in $\R$ such that $\tilde{\sigma}(0) = 0$ and $\pi \circ \tilde{\sigma} = \sigma$.
\end{lemma}
\begin{proof}{Proof:}
    Let $U = S^1 - \{-1\}$ and $V = S^1 - \{1\}$. Consider
    \[ \pi^{-1}(U) = \bigcup_{n \in \Z} \left(n - \frac{1}{2}, n + \frac{1}{2}\right) \quand \pi^{-1}(V) = \bigcup_{n \in \Z} (n, n + 1) \subseteq \R \]
    $U$ and $V$ are an open cover of $S^1$, so $\sigma^{-1}(U)$ and $\sigma^{-1}(V)$ are an open cover of $[0, 1]$. By the Lebesgue lemma, find an $\ell > 0$ such that every interval $A$ with $\abs{A} < \ell$ lies in $\sigma^{-1}(U)$ or $\sigma^{-1}(V)$, so $\sigma(A) \subseteq U$ or $\sigma(A) \subseteq V$.\\

    Let $0 = T_0 < T_1 < \cdots < T_m = 1$ so that $\abs{T_{j + 1} - T_j} < \ell$ for all $j$. Since $\sigma(0) = 1 \in U$ and $\sigma(0) \notin V$, we have $\sigma[T_0, T_1] \subseteq U$. Then, as $\pi\res{(-\frac{1}{2}, \frac{1}{2})} : (-\frac{1}{2}, \frac{1}{2}) \to U$ is a homeomorphism, let $f$ be the inverse. Then, define
    \[ \tilde{\sigma}(s) = f(\sigma(s)) \tt{ for } s \in [T_0, T_1] \]
    Inductively, suppose $\tilde{\sigma}$ is defined for $s \in [0, T_k]$. We want to define $\tilde{\sigma}$ on $[T_k, T_{k + 1}]$.\\

    Suppose $\sigma([T_k, T_{k + 1}]) \subseteq U$ and $\tilde{\sigma}(T_k) \in (n - \frac{1}{2}, n + \frac{1}{2})$. Now, $\pi\res{(n - \frac{1}{2}, n + \frac{1}{2})} : (n - \frac{1}{2}, n + \frac{1}{2}) \to U$ is a homeomorphism, so say $g$ is its inverse. Then, define
    \[ \tilde{\sigma}(s) = g(\sigma(s)) \tt{ for } s \in [T_k, T_{k + 1}] \]
    so $\tilde{\sigma}$ is defined on $[0, T_{k + 1}]$.\\

    Now, suppose $\sigma([T_k, T_{k + 1}]) \subseteq V$ and $\tilde{\sigma}(T_k) \in (n, n + 1)$. Let $h$ be the inverse of $\pi\res{(n, n + 1)} : (n, n + 1) \to V$ and define
    \[ \tilde{\sigma}(s) = h(\sigma(s)) \tt{ for } s \in [T_k, T_{k + 1}] \]
    so $\tilde{\sigma}$ is defined on $[0, T_{k + 1}]$. Repeat until we lift all of $\sigma$.
    \fig{0.5}{Lecture19Fig2}{}

    $\qed$
\end{proof}

\begin{lemma}[\textbf{Lemma 19.2:} Homotopy Lifting Lemma]
    If $F : I \times I \to S^1$ is such that $F(0, t) = 1 \in S^1$, then there exists a unique $\widetilde{F} : I \times I \to \R$ such that $\widetilde{F}(0, t) = 0$ for all $t \in [0, 1]$ and $F = \pi \circ \widetilde{F}$.
\end{lemma}
\begin{proof}{Proof:}
    Subdivide $[0, 1] \times [0, 1]$ into squares so $F$ sends each square to $U$ or $V$ (Lebesgue lemma), and lift the squares starting with the lower left.
    \fig{0.6}{Lecture19Fig3}{}
\end{proof}

\begin{theorem}[\textbf{Theorem 19.3:} Fundamental Group of the Circle]
    The function $\phi : \Z \to \pi_1(S^1, 1)$ given by $\phi(n) = \brac{\pi \circ \gamma_n}$ is an isomorphism.
\end{theorem}
\begin{proof}{Proof:}
    First, we show $\phi$ is a homomorphism. Given $m, n \in \Z$, say $\sigma(s) = \gamma_n(s) + m$ for $0 \leq s \leq 1$. Then, $\pi \circ \sigma = \pi \circ \gamma_n$, and $\gamma_m \cdot \sigma$ joins $0$ to $m + n$, so 
    \begin{align*}
        \phi(m + n) = \brac{\pi \circ \gamma_{m + n}} &= \brac{\pi \circ (\gamma_m \cdot \sigma)} = \brac{\pi \circ \gamma_m} \cdot \brac{\pi \circ \sigma}\\ 
        &= \brac{\pi \circ \gamma_m} \cdot \brac{\pi \circ \gamma_n}\\
        &= \phi(m) \cdot \phi(n)
    \end{align*}

    Now, we need to show $\phi$ is onto. Given $\brac{\alpha} \in \pi_1(S^1, 1)$, lift $\alpha$ to a path $\tilde{\alpha} : [0, 1] \to \R$. Since $\alpha(1) = 1$, $\tilde{\alpha}(1) \in \Z$. Say $\tilde{\alpha}(1) = n$. Then, by convexity, $\tilde{\alpha} \simeq \gamma_n$ (rel $\{0, 1\}$). Thus, $\alpha = \pi \circ \tilde{\alpha} \simeq \pi \circ \gamma_n$, so
    \[ \brac{\alpha} = \brac{\pi \circ \gamma_n} = \phi(n) \]

    Now, we show $\phi$ is one-to-one. Suppose $\phi(n) = \brac{e}$. Then $\pi \circ \gamma_n \simeq e$. Say $F : I \times I \to S^1$ is the homotopy from $e$ to $\pi \circ \gamma_n$.
    \fig{0.6}{Lecture19Fig4}{}

    By the homotopy lifting lemma, lift $F$ to $\widetilde{F} : [0, 1] \times [0, 1] \to \R$. Then, $\widetilde{F}(0, t) = 0$ for all $t \in [0, 1]$.\\

    Now, let $P = \{\tt{left edge}\} \cup \{\tt{right edge}\} \cup \{\tt{bottom edge}\}$. Now, $\pi \circ \widetilde{F}(P) = F(P) = 1$ and $P$ is connected, so $\widetilde{F}(P) = N$ for a single integer $N$ (otherwise, it must go to another integer, contradicting connectedness).\\

    However, $\widetilde{F}(0, 0) = 0$, so $N = 0$, so $\widetilde{F}\res{P} = 0$. Now, on the top edge, $\{\widetilde{F}(s, 1)\}_{s \in [0, 1]}$ is a lift of $\pi \circ \gamma_n$ to $\R$, so $\widetilde{F}(s, 1) = \gamma_n(s) + m$ for some $m \in \Z$. Plugging in $s = 0$, we get $m = 0$.\\

    Thus, $\widetilde{F}(s, 1) = \gamma_n(s)$ for all $s \in [0, 1]$, so $\widetilde{F}(1, 1) = \gamma_n(1)$. But $\gamma_n(1) = n$, and as $\widetilde{F}(1, 1) = 0$, $n = 0$. Thus, $\phi$ is one-to-one.

    $\qed$
\end{proof}
